\section{Who closes the road: forward deployed engineers}
\label{sec:fde}

Like Section~\ref{sec:implications}, this section states implications of the thesis: the author's
expectation, not a result. The law says what becomes software. It does not happen by itself. In practice the work is done by a
particular kind of engineer. A forward deployed engineer (FDE) is a software engineer embedded with a
customer who designs, builds, and ships a solution inside the customer's real systems and owns the
outcome, not only the delivery. The role was coined at Palantir and has spread across the AI industry
because the hard part of enterprise AI is deployment, not the model \citep{iit2026fde}.

An FDE meets the DGF twice. First as an obstacle: an agent project passes through the same IT,
Architecture, Security, Readiness, and General gates as any other change, and the Build case of
Section~\ref{sec:routes-worked} is exactly such a project. Then as an object: the FDE is the person
who turns a general agent platform into gate implementations. That work is concrete. It means reading
the policies as a specialist would; connecting the evidence sources, from repositories and identity
systems to contract stores, ticketing, and monitoring; encoding each contract and its mandate with the
authorized principal; building adjudicated evaluation sets with the specialists; keeping the producing
agent, the verifying agent, and the authorizing mechanism separate; and instrumenting the labor
account of Section~\ref{sec:frontier}, so that the result is measured in human hours rather than in
demonstrations.

\textbf{The best FDEs will design the automation of the whole DGF, not of a single gate.} The difference is
not effort, and it is not the order of deployment; it is what gets designed. A coherent end-to-end
architecture can be deployed gate by gate: a first gate can run autonomously while using the data
model, the identities, and the interfaces designed for the next ones. What must be avoided are
automations designed in isolation: local copilots whose outputs are re-extracted and re-verified by the
next team, incompatible records, and local savings obtained by moving work elsewhere.
Section~\ref{sec:example} shows that the total human workload can then rise. An end-to-end design
builds one evidence record that every role reads and extends, because the
artifact of one role is the document of the next (Section~\ref{sec:roles}); one case state whose
conditions are carried from gate to gate; the route orders the enterprise has chosen, preserved; mandates negotiated once
per class of decision rather than once per tool; one assurance architecture; and one labor ledger
across the nine populations. The economic content of that difference is the interaction term of
Equation~\eqref{eq:interaction}, which can be positive or negative: the value of designing across gates
is the part of the saving that no sum of isolated projects delivers, and it has to be measured, not
attributed in advance to the engineer's profile. The second experiment of Section~\ref{sec:testing}
measures it.

Three consequences follow. An FDE hired by one department tends to automate that department's gate and
to stop at assistance, because the mandate, the evidence, and the savings that would justify going
further belong to other departments; end-to-end automation needs a sponsor above the gates, typically
the owner of the DGF itself. The patterns an FDE builds for Buy, Integrate, and Build are reusable across
customers, which is the shared-provision mechanism of Section~\ref{sec:provision} in human form: the
first deployments are expensive and later ones amortize them. Reuse is one way a low-volume route such
as Integrate becomes economical; lowering its exception effort, its review, or its upkeep are others. And the FDE's own work belongs in the account: it is a transition
investment while the road is being built and recurring support $B_A$ while it must be maintained.
Law~5 applies to it as to any other supervisory work.

The profile that matters is therefore not the engineer who can make one agent pass one review. It is
the engineer who can read a DGF as a chain of documents and artifacts, redesign the handoffs, obtain
the mandates, and prove with the labor account that human work fell across the whole chain. This is
the author's assessment of where the scarce skill will lie. It can be checked by comparing, at
comparable scope and conditions, the residual ratio $\rho$ of programs designed with common interfaces
and one labor account with that of programs designed in silos; deployment can proceed gate by gate in
both groups.

\subsection{After the DGF: the professions of the business itself}
\label{sec:afterdgf}

The DGF gathers professions of information technology and its governance that exist in every
enterprise, whatever it sells. That is why it comes first in the comparative hypothesis of
Section~\ref{sec:candidate}, and why it is the natural first road for an FDE: a pattern built for one
enterprise's gates is reusable in every other enterprise. The thesis of this section goes one step
further. \textbf{After the DGF, it will be the turn of the professions that are the business itself:
the underwriter and the claims adjuster in insurance, the loan officer and the credit analyst in
banking, and so on in each industry.} These professions will be transformed, and many of their posts
removed, by the same mechanism, because their work has the same shape: a submission, a claim file, or a
loan application is read, and a priced quote, a coverage decision, or a credit decision is returned.

The movement has already begun on the standard path of these professions. In insurance claims,
Allianz's Project Nemo coordinates seven specialized agents that handle coverage checks, weather
verification, fraud screening, and payout calculation for low-value food-spoilage claims in Australia;
the insurer reports processing and settlement times cut by about 80\%, from days to hours, with a claims
professional still making the final payout decision \citep{allianz2025}. That figure is a delay, not a
count of human hours; and because a person confirms every payout, the remaining human work sits in the
review term of the account rather than in a small exception queue. In commercial underwriting, AIG has built an
agentic system with Palantir and Anthropic: its chief executive described reviewing every submission
in some lines without adding underwriters, then a multi-agent design in which agents ingest
submissions, evaluate risks against underwriting guidelines, benchmark pricing, and synthesize their
findings for an underwriter, who keeps the decision
\citep{insurancejournal2025,reinsurancenews2026}. In consumer lending, Upstart reports that more than 90\% of its funded loans
are fully automated, with no human involvement required by the company, while its presentation for the
same quarter shows about a quarter of evaluated applications going to manual review
\citep{upstart2026}. The two figures describe different populations with different conversion rates,
and the company's definition of end to end stops at funding for personal loans but at approval for auto
loans and home-equity lines, where human action is still required before funding. It is a textbook
case of why the accounting boundary of Section~\ref{sec:frontier} matters.

These cases are consistent with the patterns of Section~\ref{sec:hypotheses}; they do not measure them.
The standard path goes first and a residual remains, but H3 would need the share of baseline labor held
by that residual, which none of these sources reports. Analysis is delegated while authorization is
retained, the adjuster still releasing the payment and the underwriter still binding the risk, which is
compatible with H5 without measuring the two trajectories. None of them establishes H6. And the vendors that deploy these systems work
through embedded engineers; the FDE model itself comes from one of AIG's partners. What is automated
today is the simple claim and the small personal loan, not the complex commercial risk, the disputed
claim, or the corporate credit. The law predicts that those follow, by the same exhaustion of the
residual (H6).

Two differences explain the order. These decisions often concern persons, and the law treats them
accordingly: creditworthiness assessment and risk pricing in life and health insurance are listed as
high-risk uses in the AI Act \citep[Annex~III]{aiact}, and solely automated decisions about individuals
fall under GDPR Article~22 \citep{gdpr}. And these professions are vertical: an underwriting
pattern is reusable across insurers, not across all enterprises, so the amortization base of
Section~\ref{sec:provision} is narrower than for the DGF. Against that, the volumes are far larger than
a governance pool of 140 people, so the labor at stake is far greater. That the standard path of these
professions is advancing in parallel is also the most direct challenge to the ranking of
Section~\ref{sec:candidate}, which remains a hypothesis.

The two movements are linked. Every automation of a business profession is itself a project that must
pass the enterprise's gates, as the Build case shows: an automated DGF removes the brake on all the
automation that follows it. And the FDE who has closed a whole DGF has learned the method that the
business chains require: contracts for each step, one evidence record, mandates by class of decision,
an assurance architecture, and a labor account that proves the result. Submission, pricing, binding,
and claims in insurance, or application, identity checks, credit decision, and servicing in banking,
are roads of the same kind. The same reading applies to a mortgage processor, a customs declarant, or a
pharmacovigilance case handler. The labor
account of Section~\ref{sec:frontier} is the instrument that will confirm or refute this expectation,
profession by profession.
